\chapter{Related Work}
\section{No-Reference CT Image Quality Assessment}
Traditional full-reference IQA metrics such as PSNR and SSIM \cite{wang2004image} quantify deviations from a reference image, while learned perceptual distances such as LPIPS \cite{zhang2018unreasonable} compare deep feature representations. Statistical image-information criteria have also been used to connect visual fidelity with natural-image regularities \cite{sheikh2006statistical}. Although useful for reconstruction evaluation, these generic fidelity measures do not directly model radiologists' judgments and often require a suitable reference image. No-reference IQA instead predicts perceptual quality from the observed image alone. For CT imaging, SSIQA introduced self-supervised noise-level prediction as an auxiliary task for no-reference quality assessment \cite{imran2021ssiqa}. The LDCTIQAC 2023 benchmark subsequently enabled direct comparison of CT-IQA methods using radiologists' quality scores for images affected by noise and sparse-view streak artifacts \cite{lee2025low}.

Recent CT-specific methods increasingly rely on learned representations. Shi et al.\ estimate primary image content with a denoising diffusion model and combine a dissimilarity map with a transformer-based evaluator \cite{shi2024blind}. DistilIQA uses convolutional and transformer layers together with teacher--student distillation \cite{baldeoncalisto2024distiliqa}. TFKT V2 transfers knowledge from MOS-annotated natural-image datasets to CT IQA through a hybrid convolutional--transformer model \cite{rifa2025tfkt}. These approaches predict quality without requiring a high-dose reference image, but they remain specialized IQA architectures.

\section{Prompt-Driven and Vision-Language IQA}
Multimodal large language models offer a more general route to image quality assessment. Wu et al.\ systematically evaluated prompting strategies for IQA and showed that general-purpose multimodal models can be used as quality assessors \cite{wu2024comprehensive}. In medical imaging, MedIQA introduces a prompt-driven foundation model that adapts across modalities and clinical scenarios \cite{xun2025mediqa}. More recently, CAP-IQA combines radiology-style prompts with context-aware fusion and causal debiasing for abdominal CT quality assessment \cite{rifa2026capiqa}. In contrast to these systems, we adapt an open medical VLM, MedGemma 1.5 \cite{sellergren2026medgemma}, to generate radiologist-aligned LDCT MOS predictions directly.

\section{Parameter-Efficient and Reward-Based Adaptation}
Parameter-efficient fine-tuning is particularly relevant when adapting large models with limited domain-specific data. LoRA introduces trainable low-rank updates while keeping the pretrained model weights frozen \cite{hu2022lora}. Beyond supervised adaptation, GRPO provides a memory-efficient reinforcement learning objective that optimizes groups of generated responses according to a task reward \cite{shao2024deepseekmath}. Zhu et al.\ investigate GRPO-based fine-tuning for medical visual question answering and study factors such as model initialization and medical semantic alignment \cite{zhu2025medicalgrpo}. Our work examines a complementary scalar-prediction setting, where the target is calibrated LDCT MOS rather than a reasoning response.
